\section{NPT 的范式转换：预测时显式读取训练数据}

现在把 attention 的对象从 tokens 换成 datapoints。标准 parametric model 在训练后把经验压进参数 $\theta$，推理一条样本时通常只读取该样本；NPT 则让训练点与测试点共同进入模型，使预测对当前可见数据集保持显式依赖。

\lecturefigure{slide-10-npt-overview.jpg}{NPT 一页概览：dataset-as-input、datapoint self-attention 与 stochastic masking。}{Stanford Online 官方视频 00:17:00--00:17:21。}

这张图给出三个不可删减的组成部分：第一，输入不是单行而是数据集；第二，模型在 rows 之间做 self-attention；第三，masking objective 决定哪些值可见、哪些值需要重建。只保留其中任意两项都不能得到论文所称的 NPT predictive mechanism。

\lecturefigure{slide-11-npt-motivation.jpg}{Parametric prediction 与 NPT：后者让目标行显式读取其他 datapoints。}{Stanford Online 官方视频 00:17:22--00:18:17。}

\begin{importantbox}{本讲中“non-parametric”的精确定义}
Non-parametric 不等于“没有参数”。NPT 仍有大量可学习的 attention、embedding 与 feed-forward 参数。这里强调的是：预测函数在 inference 时显式依赖 supplied training data，写作 $p(y_*\mid x_*,\mathcal{D}_{\mathrm{train}})$，而不是只依赖 $x_*$ 与固定参数 $\theta$。
\end{importantbox}

\teachervoice{有人在社交媒体上把 NPT 称作 “k-NN 2.0”。这个说法适合建立直觉：模型会从其他 rows 查找信息；但 NPT 不只按固定距离找邻居，它可以学习多层、任务相关、非线性的 match--transform--aggregate 过程。}

\subsection{传统 non-parametric models 在哪里}

NPT 并未发明 non-parametric prediction。Gaussian process、kernel method 与 k-nearest neighbors 都让预测显式依赖训练数据；deep kernel learning 与 neural process 等工作也尝试结合 representation learning。NPT 的主张更窄：Transformer 可以作为一个通用关系学习器，把检索、比较与聚合共同端到端训练。

\lecturefigure{slide-12-traditional-nonparametric.jpg}{传统 non-parametric models 与深度扩展：NPT 的历史位置。}{Stanford Online 官方视频 00:22:00--00:23:39。}

\begin{knowledgebox}{术语消化：几类 explicit-data prediction}
\begin{tabularx}{\textwidth}{L{0.20\textwidth}L{0.27\textwidth}X}
\toprule
方法 & 如何使用训练数据 & 与 NPT 的关系 \\
\midrule
k-NN & 按预设距离检索邻居并投票/平均 & NPT 可学习相似度、聚合与后续变换，不局限于固定距离。 \\
Kernel method & 通过 kernel similarity 加权训练点 & Attention 同样形成数据依赖权重，但具有深层表示和多层组合。 \\
Gaussian process & kernel 定义函数分布，并给出预测不确定性 & NPT 不是标准 GP，也不自动具有一致的 Bayesian uncertainty。 \\
Deep kernel / neural process & 用神经网络学习表示或条件聚合 & 与 NPT 共享“参数化关系学习 + 显式 context”的动机。 \\
\bottomrule
\end{tabularx}
\end{knowledgebox}

\teachervoice{Kossen 强调他们想要的是易用、plug-and-play、能适应多种场景的 non-parametric mechanism，而不是依赖复杂专用 inference scheme。这个目标是设计动机，不应被误读成已经证明 NPT 在所有场景都优于经典方法。}

\begin{warningbox}{NPT 不是本讲已经完成的 meta-learning 系统}
论文实验主要为一个固定数据集训练一个模型，并假设 train/test 来自同一任务分布。Meta-learning 通常要求跨多个任务或数据集迁移。NPT 可以被扩展为把新 dataset 当 context 的系统，但那是讲者讨论的未来方向，不是本讲实验已经证明的能力。
\end{warningbox}

\teachervoice{Q\&A 中三位讲者把边界说得很清楚：当前工作仍是标准 supervised learning；多数据集、zero-shot 或少量 gradient step 的 meta-learning 只是值得探索的 framing。}

\subsection{本章小结}

NPT 的新意不是“无参数”，而是把训练 datapoints 保留在预测计算图里，并让 attention 学习如何读取它们。它继承经典 non-parametric intuition，同时以深度表示与多层关系计算扩大可学习机制的范围。

